\documentclass[12pt]{article}
\usepackage[utf8]{inputenc}
\usepackage[T1]{fontenc}
\usepackage{hyperref}
\title{Kimki Language Resources for K-12 Curriculum Development}
\author{Compiled by Claude Sonnet 5.5 (claude-sonnet-5-5) for Andr\'as Kornai}
\date{}
\begin{document}
\maketitle

\section{Introduction}
Kimki (ISO 639-3: sbt; also Aipki, Kimgi, Sukubatom, Sukubatong) is spoken by about 500 people (2004) in Batom District, Pegunungan Bintang Regency, Papua, Indonesia. Foley treats it as an isolate. Usher links it to South Pauwasi. Dialects are recorded from Batom and Sabi villages. Vitality sources disagree. Very little documentation exists, and no school materials were found.

\section{Primary Educational Resources}
No primer, textbook or school programme was found.

\section{Linguistic and Reference Materials}
\begin{itemize}
\item Wikipedia, Kimki language: basic vocabulary and sample sentences, citing Foley (2018). \url{https://en.wikipedia.org/wiki/Kimki_language}
\item Rumaropen (2004), draft sociolinguistic survey of Kimki varieties in the southeast of Mount Ji (Glottolog record). \url{https://glottolog.org/resource/reference/id/45210}
\item ASJP Kimki wordlist, derived from a 1980 wordlist collected by Reimer. \url{https://asjp.clld.org/languages/KIMKI}
\item Numeral data for Kimki. \url{https://lingweb.eva.mpg.de/channumerals/Kimki.htm}
\item Endangered Languages Project record. \url{https://endangeredlanguages.com/lang/5147}
\end{itemize}

\section{Religious and Community Resources}
No scripture translations or community organisations were found. Joshua Project has a people-group page (Sukubatong, Kimki). \url{https://joshuaproject.net/people_groups/15111}

\section{Digital and Audio Resources}
No audio, video or corpus resources were found beyond the wordlists above.

\section{Curriculum Development Context}
There is no known orthography, literacy material or school use. Indonesian is the likely language of instruction. Any curriculum would start with documentation work such as a dictionary and an orthography, working with the community.
\end{document}